% Bab: FAOA-2015-CH03
% Provenans: OpenAI Codex gpt-5.6-sol, Ultra, atas arahan pengguna
% Lisensi komponen solusi asli: CC BY-SA 4.0
% Non-endorsement: solusi ini bukan karya John M. Erdman dan tidak menyiratkan dukungannya.
% Jumlah solusi: 7
\markboth{SOLUSI O001 UNTUK BAB 3: RUANG LINEAR BERNORMA}{SOLUSI O001 UNTUK BAB 3: RUANG LINEAR BERNORMA}
\section*{Solusi O001 untuk Bab 3: Ruang Linear Bernorma}
\addcontentsline{toc}{section}{Solusi O001 untuk Bab 3: Ruang Linear Bernorma}
\noindent\textit{Catatan komponen.}
Ketujuh solusi berikut ditulis terpisah dari teks sumber, dilisensikan di bawah
CC BY-SA 4.0, dan berprovenans
\emph{OpenAI Codex gpt-5.6-sol, Ultra, atas arahan pengguna}.
Solusi-solusi ini bukan tulisan John M. Erdman dan tidak menyiratkan dukungannya.

% O001-SOLUTION-ID: O001-FAOA-2015-CH03-EX-001-SOLUTION
% SOURCE-EXERCISE-ID: FAOA-2015-CH03-NODE-0041
% STATEMENT-TARGET-SHA256: 1f693891073e0afb70f5764271328dd22974b4c3b33055ec494a61b7713f62e2
\begin{o001solution}
  {O001-FAOA-2015-CH03-EX-001-SOLUTION}
  {FAOA-2015-CH03-NODE-0041}
  {1f693891073e0afb70f5764271328dd22974b4c3b33055ec494a61b7713f62e2}
\begin{o001statement}
\label{C023191} Misalkan
 \index{l@$l_1$!sebagai ruang linear bernorma}%
 \index{l@$l_1$!barisan yang dapat dijumlahkan mutlak}%
 \index{bernorma!ruang linear!$l_1$ sebagai}%
$l_1$ merupakan himpunan semua barisan $x = (x_n)$ dari bilangan kompleks sedemikian
sehingga deret $\sum x_k$ konvergen mutlak. Jadikan $l_1$ ruang vektor dengan definisi
penjumlahan dan perkalian skalar titik demi titik yang biasa. Untuk setiap $x \in l_1$,
definisikan
 \index{<unopn@$\norm{\hphantom{x}}_1$ (norma-$1$)}%
 \index{<unopn@$\norm{\hphantom{x}}_\infty$ (norma seragam)}%
   \begin{align*}
      \norm{x}_1 &:= \sum_{k=1}^\infty \abs{x_k} \\
   \intertext{dan}
      \norm{x}_\infty &:= \sup\{\abs{x_k}\colon k \in \N\}.
   \end{align*}
(Masing-masing adalah
 \index{satu@norma-$1$!dari suatu barisan}%
 \index{norma!$1$- (dari suatu barisan)}%
norma-$1$ dan
 \index{seragam!norma!dari suatu barisan}%
\emph{norma seragam}.) Tunjukkan bahwa $\norm{\hphantom{x}}_1$ dan
$\norm{\hphantom{x}}_\infty$ keduanya merupakan norma pada~$l_1$. Kemudian buktikan atau
sangkal:
 \begin{enumerate}
  \item[(a)] Jika suatu barisan $(x^i)$ dari vektor-vektor dalam ruang linear bernorma
$\bigl(l_1, \norm{\hphantom{x}}_1\bigr)$ konvergen, maka barisan tersebut juga konvergen
dalam $\bigl(l_1, \norm{\hphantom{x}}_\infty\bigr)$.
  \item[(b)] Jika suatu barisan $(x^i)$ dari vektor-vektor dalam ruang linear bernorma
$\bigl(l_1, \norm{\hphantom{x}}_\infty\bigr)$ konvergen, maka barisan tersebut juga
konvergen dalam $\bigl(l_1, \norm{\hphantom{x}}_1\bigr)$.
 \end{enumerate}
\end{o001statement}
\begin{o001answer}
Kedua fungsi tersebut adalah norma pada $l_1$, dan
$\norm{x}_\infty\le\norm{x}_1$ untuk setiap $x\in l_1$. Karena itu (a) benar.
Pernyataan (b) salah: barisan
\[
x^i=(\underbrace{1/i,\ldots,1/i}_{i\ \text{suku}},0,0,\ldots)
\]
konvergen ke nol dalam norma seragam, tetapi
$\norm{x^i}_1=1$ untuk setiap $i$.
\end{o001answer}
\begin{o001proof}
Untuk $x\in l_1$, setiap $|x_k|$ tidak melebihi
$\sum_{j=1}^\infty|x_j|$, sehingga supremumnya berhingga. Jadi kedua rumus
memang bernilai real berhingga pada seluruh $l_1$.

Untuk norma-$1$, ketaknegatifan dan homogenitas mengikuti dari sifat nilai
mutlak. Jika $\norm{x}_1=0$, setiap suku $|x_k|$ bernilai nol, sehingga
$x=0$. Ketaksamaan segitiga titik demi titik dan penjumlahan deret
nonnegatif memberi
\[
\norm{x+y}_1
=\sum_{k=1}^\infty|x_k+y_k|
\le\sum_{k=1}^\infty(|x_k|+|y_k|)
=\norm{x}_1+\norm{y}_1.
\]
Maka $\norm{\hphantom{x}}_1$ adalah norma.

Untuk norma seragam, ketaknegatifan dan homogenitas juga mengikuti dari nilai
mutlak. Jika $\norm{x}_\infty=0$, maka $|x_k|=0$ untuk setiap $k$, sehingga
$x=0$. Selain itu,
\[
|x_k+y_k|\le|x_k|+|y_k|
\le\norm{x}_\infty+\norm{y}_\infty
\]
untuk setiap $k$; pengambilan supremum membuktikan ketaksamaan segitiga.
Jadi $\norm{\hphantom{x}}_\infty$ juga norma.

Karena
\[
\norm{z}_\infty=\sup_k|z_k|
\le\sum_{k=1}^\infty|z_k|=\norm{z}_1,
\]
konvergensi $x^i\to x$ dalam norma-$1$ mengakibatkan
$\norm{x^i-x}_\infty\le\norm{x^i-x}_1\to0$. Ini membuktikan (a).

Untuk vektor $x^i$ dalam jawaban, setiap vektor bertumpuan hingga, jadi
$x^i\in l_1$. Akan tetapi,
\[
\norm{x^i}_\infty=\frac1i\longrightarrow0,
\qquad
\norm{x^i}_1=i\frac1i=1.
\]
Dengan demikian $x^i\to0$ dalam norma seragam tetapi tidak dalam norma-$1$,
yang menyangkal (b).
\end{o001proof}
\end{o001solution}

% O001-SOLUTION-ID: O001-FAOA-2015-CH03-EX-002-SOLUTION
% SOURCE-EXERCISE-ID: FAOA-2015-CH03-NODE-0058
% STATEMENT-TARGET-SHA256: c063bcd29e442aa5670f59b396568ce8af222f0543fab7295967f54568271718
\begin{o001solution}
  {O001-FAOA-2015-CH03-EX-002-SOLUTION}
  {FAOA-2015-CH03-NODE-0058}
  {c063bcd29e442aa5670f59b396568ce8af222f0543fab7295967f54568271718}
\begin{o001statement}
Periksalah pernyataan-pernyataan yang belum dibuktikan dalam paragraf
sebelumnya. Secara khusus, buktikan bahwa
 \begin{enumerate}
    \item[(a)] $\cat{NLS_{\boldsymbol\infty}}$ merupakan suatu kategori.
    \item[(b)] Suatu isomorfisme dalam $\cat{NLS_{\boldsymbol\infty}}$ merupakan
isomorfisme ruang vektor sekaligus homeomorfisme.
    \item[(c)] Suatu pemetaan linear antara ruang-ruang linear bernorma merupakan isometri
jika dan hanya jika pemetaan tersebut mempertahankan norma.
    \item[(d)] $\cat{NLS_{\vc 1}}$ merupakan suatu kategori.
    \item[(e)] Suatu isomorfisme dalam $\cat{NLS_{\vc 1}}$ merupakan isometri sekaligus
isomorfisme ruang vektor.
    \item[(f)] Kategori $\cat{NLS_{\vc 1}}$ merupakan subkategori dari
$\cat{NLS_{\boldsymbol\infty}}$ dalam arti bahwa setiap morfisme kategori pertama
termasuk dalam kategori kedua.
 \end{enumerate}
\end{o001statement}
\begin{o001answer}
Pemetaan linear kontinu tertutup terhadap identitas dan komposisi, demikian
pula kontraksi linear. Isomorfisme dalam
$\cat{NLS_{\boldsymbol\infty}}$ adalah isomorfisme linear yang pemetaan dan
inversnya kontinu. Isomorfisme dalam $\cat{NLS_{\vc1}}$ adalah isometri
linear bijektif. Untuk pemetaan linear, sifat isometri ekuivalen dengan
pelestarian norma, dan setiap kontraksi linear merupakan pemetaan linear
kontinu.
\end{o001answer}
\begin{o001proof}
\begin{enumerate}
\item[(a)] Untuk setiap ruang linear bernorma $V$, identitas $I_V$ linear dan
kontinu. Jika $S\colon U\to V$ dan $T\colon V\to W$ linear kontinu, maka
$TS$ linear kontinu; dalam bentuk norma operator,
$\norm{TS}\le\norm{T}\norm{S}$. Asosiativitas komposisi dan hukum identitas
diwarisi dari komposisi fungsi. Semua aksioma kategori karena itu dipenuhi.

\item[(b)] Jika $T\colon V\to W$ isomorfisme dalam
$\cat{NLS_{\boldsymbol\infty}}$, terdapat morfisme
$S\colon W\to V$ dengan $ST=I_V$ dan $TS=I_W$. Maka $T$ bijektif,
$S=T^{-1}$, dan keduanya linear, jadi $T$ isomorfisme ruang vektor.
Karena $T$ dan $T^{-1}$ kontinu, $T$ juga homeomorfisme.

\item[(c)] Jika $T$ linear dan isometri, maka $T0=0$ dan
\[
\norm{Tx}=\norm{Tx-T0}=\norm{x-0}=\norm{x}.
\]
Sebaliknya, jika $T$ mempertahankan norma, linearitas memberi
\[
\norm{Tx-Ty}=\norm{T(x-y)}=\norm{x-y}
\]
untuk setiap $x,y$, sehingga $T$ isometri.

\item[(d)] Identitas memenuhi $\norm{I_Vx}=\norm{x}$, jadi merupakan
kontraksi. Jika $S$ dan $T$ kontraksi, maka
\[
\norm{TSx}\le\norm{Sx}\le\norm{x},
\]
sehingga komposisinya kontraksi. Asosiativitas dan hukum identitas kembali
diwarisi dari fungsi. Maka $\cat{NLS_{\vc1}}$ adalah kategori.

\item[(e)] Jika $T$ isomorfisme dalam $\cat{NLS_{\vc1}}$, baik $T$ maupun
$T^{-1}$ merupakan kontraksi. Jadi
\[
\norm{Tx}\le\norm{x}
=\norm{T^{-1}Tx}\le\norm{Tx}.
\]
Dengan demikian $\norm{Tx}=\norm{x}$ untuk setiap $x$. Bagian (c)
menunjukkan bahwa $T$ isometri; persamaan invers menunjukkan bahwa $T$
sekaligus isomorfisme ruang vektor.

\item[(f)] Jika $T$ kontraksi linear, maka
$\norm{Tx}\le\norm{x}$, jadi $T$ terbatas dengan $\norm T\le1$ dan karena
itu kontinu. Objek, identitas, dan komposisinya sama dengan yang diwarisi
dari $\cat{NLS_{\boldsymbol\infty}}$. Jadi setiap morfisme kategori
geometris merupakan morfisme kategori topologis, tepat dalam arti
subkategori yang dinyatakan dalam soal.
\end{enumerate}
\end{o001proof}
\end{o001solution}

% O001-SOLUTION-ID: O001-FAOA-2015-CH03-EX-003-SOLUTION
% SOURCE-EXERCISE-ID: FAOA-2015-CH03-NODE-0073
% STATEMENT-TARGET-SHA256: 730caa86f3c8bec4a40678e8832cce25c3d05b021a5fe0770bcff3a73f4a3ec2
\begin{o001solution}
  {O001-FAOA-2015-CH03-EX-003-SOLUTION}
  {FAOA-2015-CH03-NODE-0073}
  {730caa86f3c8bec4a40678e8832cce25c3d05b021a5fe0770bcff3a73f4a3ec2}
\begin{o001statement}
Verifikasilah pernyataan-pernyataan dalam Definisi~\ref{quo002def}. Secara khusus,
tunjukkan bahwa $\sim$ merupakan relasi ekuivalensi, bahwa penjumlahan dan perkalian skalar pada
himpunan kelas-kelas ekuivalensi terdefinisi dengan baik, bahwa terhadap operasi-operasi ini $V/M$ merupakan
ruang vektor, bahwa fungsi yang di sini disebut ``pemetaan hasil bagi'' memang merupakan pemetaan hasil bagi
dalam pengertian~\ref{quo001def}, dan bahwa pemetaan hasil bagi ini bersifat linear.
\end{o001statement}
\begin{o001answer}
Relasi $x\sim y\iff y-x\in M$ adalah relasi ekuivalensi. Rumus
$[x]+[y]=[x+y]$ dan $\alpha[x]=[\alpha x]$ tidak bergantung pada wakil,
menjadikan $V/M$ ruang vektor, dan pemetaan
$\pi(x)=[x]$ adalah pemetaan linear surjektif. Lebih lanjut, setiap fungsi
$h\colon V/M\to W$ yang membuat $h\pi$ linear harus linear; jadi $\pi$
adalah pemetaan hasil bagi dalam pengertian~\ref{quo001def}.
\end{o001answer}
\begin{o001proof}
Karena $M$ subruang, $0\in M$, tertutup terhadap invers aditif, dan tertutup
terhadap penjumlahan. Maka $x-x=0\in M$, sehingga $x\sim x$. Jika
$x\sim y$, maka $y-x\in M$ dan $x-y=-(y-x)\in M$, sehingga $y\sim x$.
Jika $x\sim y$ dan $y\sim z$, maka
$z-x=(z-y)+(y-x)\in M$, sehingga $x\sim z$. Jadi $\sim$ relasi
ekuivalensi.

Andaikan $[x]=[x']$ dan $[y]=[y']$. Berarti $x'-x,y'-y\in M$, sehingga
\[
(x'+y')-(x+y)=(x'-x)+(y'-y)\in M.
\]
Jadi $[x'+y']=[x+y]$. Untuk $\alpha\in\K$,
\[
\alpha x'-\alpha x=\alpha(x'-x)\in M,
\]
sehingga $[\alpha x']=[\alpha x]$. Kedua operasi dengan demikian
terdefinisi dengan baik.

Setiap aksioma ruang vektor sekarang dapat diperiksa pada wakilnya. Sebagai
contoh,
\[
([x]+[y])+[z]=[(x+y)+z]=[x+(y+z)]=[x]+([y]+[z]),
\]
$[0]$ adalah identitas aditif, $[-x]$ invers dari $[x]$, dan
\[
\alpha([x]+[y])=[\alpha(x+y)]
=[\alpha x]+[\alpha y].
\]
Komutativitas, aksioma skalar lainnya, dan identitas skalar mengikuti dengan
perhitungan wakil yang sama. Maka $V/M$ adalah ruang vektor.

Pemetaan $\pi\colon V\to V/M$, $\pi(x)=[x]$, surjektif karena setiap kelas
memiliki wakil. Selain itu,
\[
\pi(x+y)=[x+y]=[x]+[y],\qquad
\pi(\alpha x)=[\alpha x]=\alpha[x],
\]
jadi $\pi$ linear.

Terakhir, ambil sebarang ruang vektor $W$ dan fungsi
$h\colon V/M\to W$ sedemikian sehingga $h\circ\pi$ linear. Untuk semua
$[x],[y]\in V/M$ dan $\alpha\in\K$,
\[
\begin{aligned}
h([x]+[y])
 &=h(\pi(x+y))
  =(h\pi)(x+y)
  =(h\pi)(x)+(h\pi)(y)
  =h([x])+h([y]),\\
h(\alpha[x])
 &=h(\pi(\alpha x))
  =(h\pi)(\alpha x)
  =\alpha(h\pi)(x)
  =\alpha h([x]).
\end{aligned}
\]
Jadi $h$ linear. Ini persis sifat universal pada
Definisi~\ref{quo001def}, sehingga $\pi$ memang pemetaan hasil bagi.
\end{o001proof}
\end{o001solution}

% O001-SOLUTION-ID: O001-FAOA-2015-CH03-EX-004-SOLUTION
% SOURCE-EXERCISE-ID: FAOA-2015-CH03-NODE-0094
% STATEMENT-TARGET-SHA256: c2de90c5ffcfd0c7af44f29f75839337bc75b7c09430a8da07d70509a7db2ad1
\begin{o001solution}
  {O001-FAOA-2015-CH03-EX-004-SOLUTION}
  {FAOA-2015-CH03-NODE-0094}
  {c2de90c5ffcfd0c7af44f29f75839337bc75b7c09430a8da07d70509a7db2ad1}
\begin{o001statement}
Berikan bukti yang \emph{sangat} singkat (tanpa $\epsilon$, $\delta$, ataupun himpunan terbuka) bahwa jika
$(x_n)$ dan $(y_n)$ merupakan barisan-barisan dalam suatu ruang linear bernorma yang masing-masing konvergen ke $a$ dan $b$,
maka $x_n + y_n \sto a + b$.
\end{o001statement}
\begin{o001answer}
\[
\norm{(x_n+y_n)-(a+b)}
\le\norm{x_n-a}+\norm{y_n-b}\longrightarrow0.
\]
Jadi $x_n+y_n\to a+b$.
\end{o001answer}
\begin{o001proof}
Dengan mengelompokkan suku dan memakai ketaksamaan segitiga,
\[
\norm{(x_n+y_n)-(a+b)}
=\norm{(x_n-a)+(y_n-b)}
\le\norm{x_n-a}+\norm{y_n-b}.
\]
Kedua suku di ruas kanan menuju nol berdasarkan kedua hipotesis, sehingga
ruas kiri juga menuju nol. Ini adalah definisi konvergensi dalam norma dan
tidak memakai $\epsilon$, $\delta$, ataupun himpunan terbuka.
\end{o001proof}
\end{o001solution}

% O001-SOLUTION-ID: O001-FAOA-2015-CH03-EX-005-SOLUTION
% SOURCE-EXERCISE-ID: FAOA-2015-CH03-NODE-0118
% STATEMENT-TARGET-SHA256: a54a95af268dbf3cafc70d6ee756695be4fe8a45d3f1d175b15012cfa280f0d5
\begin{o001solution}
  {O001-FAOA-2015-CH03-EX-005-SOLUTION}
  {FAOA-2015-CH03-NODE-0118}
  {a54a95af268dbf3cafc70d6ee756695be4fe8a45d3f1d175b15012cfa280f0d5}
\begin{o001statement}
Apa saja koproduk dalam kategori $\cat{NLS_{\boldsymbol\infty}}$?
\end{o001statement}
\begin{o001answer}
Koproduk suatu keluarga ruang linear bernorma ada tepat ketika hanya
berhingga banyak anggotanya yang bukan ruang nol. Untuk
$V_1,\ldots,V_n$, satu modelnya ialah jumlah langsung aljabar
\[
\bigoplus_{k=1}^nV_k
\quad\text{dengan}\quad
\norm{(x_1,\ldots,x_n)}=\sum_{k=1}^n\norm{x_k},
\]
beserta injeksi koordinat. Koproduk keluarga kosong adalah ruang nol.
Jika terdapat tak hingga banyak objek bukan nol, koproduknya tidak ada dalam
kategori dengan semua pemetaan linear terbatas sebagai morfisme.
\end{o001answer}
\begin{o001proof}
Untuk keluarga hingga $V_1,\ldots,V_n$, injeksi koordinat
$\iota_k\colon V_k\to\bigoplus_{j=1}^nV_j$ merupakan pemetaan linear
terbatas. Jika $f_k\colon V_k\to W$ adalah pemetaan linear terbatas,
definisikan
\[
F(x_1,\ldots,x_n)=\sum_{k=1}^nf_k(x_k).
\]
Pemetaan ini satu-satunya pemetaan linear dengan $F\iota_k=f_k$, dan
\[
\norm{F(x_1,\ldots,x_n)}
\le\sum_{k=1}^n\norm{f_k}\norm{x_k}
\le\left(\max_{1\le k\le n}\norm{f_k}\right)
   \sum_{k=1}^n\norm{x_k}.
\]
Jadi $F$ terbatas dan sifat universal koproduk dipenuhi. Ruang nol tidak
mengubah konstruksi ini, sehingga keluarga dengan hanya berhingga banyak
anggota bukan nol mempunyai koproduk yang sama.

Sekarang andaikan keluarga $(V_\lambda)_{\lambda\in\Lambda}$ mempunyai tak
hingga banyak anggota bukan nol dan mempunyai koproduk hipotetis
$(C,(\iota_\lambda))$. Untuk setiap indeks $\lambda$ dengan
$V_\lambda\ne0$, gunakan keluarga morfisme ke $V_\lambda$ yang pada
$V_\lambda$ adalah identitas dan pada semua komponen lain adalah nol.
Sifat universal menghasilkan $p_\lambda\colon C\to V_\lambda$ dengan
$p_\lambda\iota_\lambda=I_{V_\lambda}$; khususnya
$\iota_\lambda\ne0$.

Pilih indeks-indeks berbeda $\lambda_1,\lambda_2,\ldots$ dengan
$V_{\lambda_n}\ne0$. Untuk setiap $n$, tetapkan
$f_{\lambda_n}=n\,\iota_{\lambda_n}\colon V_{\lambda_n}\to C$, dan tetapkan
$f_\lambda=0$ pada indeks lain. Masing-masing $f_\lambda$ adalah morfisme
terbatas. Sifat universal akan menghasilkan pemetaan linear terbatas
$F\colon C\to C$ dengan
\[
F\iota_{\lambda_n}=n\,\iota_{\lambda_n}.
\]
Ambil $x_n$ sehingga $\iota_{\lambda_n}x_n\ne0$. Maka
\[
\norm F\ge
\frac{\norm{F\iota_{\lambda_n}x_n}}
     {\norm{\iota_{\lambda_n}x_n}}
=n
\]
untuk setiap $n$, yang mustahil bagi operator terbatas. Jadi koproduk tak
ada bila tak hingga banyak komponen bukan nol.
\end{o001proof}
\end{o001solution}

% O001-SOLUTION-ID: O001-FAOA-2015-CH03-EX-006-SOLUTION
% SOURCE-EXERCISE-ID: FAOA-2015-CH03-NODE-0119
% STATEMENT-TARGET-SHA256: c18a12a34660fdaf55f44c9058084c0551cb4ad4560ec59d554f7a37201caf37
\begin{o001solution}
  {O001-FAOA-2015-CH03-EX-006-SOLUTION}
  {FAOA-2015-CH03-NODE-0119}
  {c18a12a34660fdaf55f44c9058084c0551cb4ad4560ec59d554f7a37201caf37}
\begin{o001statement}
Tentukan hasil kali dan koproduk dua ruang $V$ dan $W$ dalam kategori
$\cat{NLS_{\vc 1}}$ yang terdiri atas ruang-ruang linear bernorma dan kontraksi-kontraksi linear.
\end{o001statement}
\begin{o001answer}
Hasil kali adalah $V\times W$ dengan norma maksimum
\[
\norm{(v,w)}_{\mathrm{prod}}=\max\{\norm v,\norm w\}
\]
dan proyeksi koordinat. Koproduk adalah $V\oplus W$ dengan norma jumlah
\[
\norm{(v,w)}_{\mathrm{koprod}}=\norm v+\norm w
\]
dan injeksi koordinat.
\end{o001answer}
\begin{o001proof}
Pada norma maksimum, kedua proyeksi
$\pi_V(v,w)=v$ dan $\pi_W(v,w)=w$ adalah kontraksi. Jika
$f\colon X\to V$ dan $g\colon X\to W$ kontraksi, satu-satunya pemetaan
linear $H\colon X\to V\times W$ yang memenuhi
$\pi_VH=f$ dan $\pi_WH=g$ ialah $H(x)=(f(x),g(x))$. Pemetaan ini kontraksi
karena
\[
\norm{H(x)}_{\mathrm{prod}}
=\max\{\norm{f(x)},\norm{g(x)}\}\le\norm{x}.
\]
Jadi objek tersebut adalah hasil kali.

Pada norma jumlah, injeksi
$\iota_V(v)=(v,0)$ dan $\iota_W(w)=(0,w)$ adalah kontraksi, bahkan isometri.
Jika $f\colon V\to X$ dan $g\colon W\to X$ kontraksi, satu-satunya pemetaan
linear $K\colon V\oplus W\to X$ dengan
$K\iota_V=f$ dan $K\iota_W=g$ ialah
$K(v,w)=f(v)+g(w)$. Ketaksamaan
\[
\norm{K(v,w)}
\le\norm{f(v)}+\norm{g(w)}
\le\norm v+\norm w
=\norm{(v,w)}_{\mathrm{koprod}}
\]
menunjukkan bahwa $K$ kontraksi. Jadi objek kedua adalah koproduk.
\end{o001proof}
\end{o001solution}

% O001-SOLUTION-ID: O001-FAOA-2015-CH03-EX-007-SOLUTION
% SOURCE-EXERCISE-ID: FAOA-2015-CH03-NODE-0163
% STATEMENT-TARGET-SHA256: 9ceb84511f36b87a9eb590f586023123c303e4bc85ba4a4938930a5c6aca9926
\begin{o001solution}
  {O001-FAOA-2015-CH03-EX-007-SOLUTION}
  {FAOA-2015-CH03-NODE-0163}
  {9ceb84511f36b87a9eb590f586023123c303e4bc85ba4a4938930a5c6aca9926}
\begin{o001statement}
Kawan Anda, Fred R. Dimm, mengira bahwa ia telah menemukan pembuktian langsung yang sangat sederhana untuk
versi teorema Hahn--Banach yang diberikan dalam Teorema~\ref{HBTIII}: Suatu fungsional
linear kontinu $f$ yang didefinisikan pada subruang (linear) $M$ dari ruang linear bernorma~$V$ dapat
diperpanjang menjadi fungsional linear kontinu $g$ pada seluruh $V$ tanpa memperbesar
normanya. Menurutnya, kita hanya perlu mengambil $N$ sebagai sebarang komplemen ruang vektor
dari~$M$ dan mengambil $B$ sebagai suatu basis (Hamel) untuk $N$. Definisikan $g(e) = 0$ untuk setiap vektor $e
\in B$ dan $g = f$ pada~$M$. Kemudian perpanjang $g$ secara linear ke seluruh~$V$. Tentunya, kata Fred, $g$ linear berdasarkan konstruksinya dan normanya tidak bertambah karena kita
telah menjadikannya nol di tempat-tempat yang sebelumnya belum didefinisikan. Tunjukkan kepada Fred letak kekeliruannya.
\end{o001statement}
\begin{o001answer}
Perpanjangan Fred adalah $g=f\circ P$, dengan $P$ proyeksi aljabar ke $M$
sepanjang komplemen $N$. Proyeksi itu tidak harus kontraktif, bahkan tidak
harus terbatas. Jadi menaruh nilai nol pada $N$ tidak mengendalikan norma
$g$. Sudah pada $\R^2$ dapat terjadi $\norm g>\norm f$.
\end{o001answer}
\begin{o001proof}
Ambil $V=\R^2$ dengan norma Euclid,
\[
M=\spn\{(1,0)\},\qquad
f(t,0)=t,\qquad
N=\spn\{(1,1)\}.
\]
Fungsional $f$ mempunyai norma $1$ pada $M$. Pilih basis Hamel
$B=\{(1,1)\}$ untuk $N$ dan lakukan persis konstruksi Fred: tetapkan
$g(1,1)=0$ dan $g=f$ pada $M$. Karena
\[
(x,y)=(x-y,0)+y(1,1),
\]
perpanjangan linearnya adalah
\[
g(x,y)=x-y.
\]
Norma fungsional ini pada ruang Euclid adalah
\[
\norm g=\sqrt{1^2+(-1)^2}=\sqrt2>1=\norm f.
\]
Sebagai pemeriksaan langsung, pada
$u=(1,-1)/\sqrt2$ berlaku $\norm u=1$ tetapi $|g(u)|=\sqrt2$.

Kesalahannya ialah bahwa vektor $m+n$ dapat mempunyai norma jauh lebih kecil
daripada norma komponen $m$. Menjadikan $g$ nol pada $n$ tidak menghilangkan
kemungkinan pembatalan geometris itu. Secara umum, jika
$P\colon M\oplus N\to M$ adalah proyeksi aljabar, konstruksi Fred memberi
$g=f\circ P$; suatu komplemen aljabar sebarang tidak menjamin
$\norm P\le1$, dan dalam ruang berdimensi tak hingga $P$ bahkan dapat tidak
terbatas. Karena itu konstruksi tersebut tidak menjamin kesamaan norma maupun
kekontinuan.
\end{o001proof}
\end{o001solution}
